\begin{figure}[ht]
\centering
\begin{tikzpicture}[
  box/.style={draw,rounded corners=2pt,inner sep=7pt,
              text width=12.5cm,align=center,font=\small},
  transfer/.style={-{Stealth[length=2mm]},thick}]
\node[box] (full) at (0,0)
  {$[R_1<\cdots<R_c=R<RT_1<\cdots<RT_{m-c}=E]$};
\node[box] (suffix) at (0,-1.45)
  {$[RT_1<\cdots<RT_{m-c}=E]$};
\node[box] (image) at (0,-2.90)
  {$[T_1<\cdots<T_{m-c}=K]$};
\draw[transfer] (full)--node[right,font=\small,fill=white,inner sep=2pt]
  {residue at $[R_1<\cdots<R_c]$} (suffix);
\draw[transfer] (suffix)--node[right,font=\small,fill=white,inner sep=2pt]
  {intersection with $H$} (image);
\end{tikzpicture}
\caption{The chosen supported flag in Lemma~\ref{lem:restriction}, with
$E=R\times K$, $K=E\cap H$, and $0<c=\rk R<m$.
Here $T_1<\cdots<T_{m-c}=K$ is a full flag of subgroups of $K$.
The maximal-index hypothesis forces every continuation vertex $B$ in
the residue support to equal $R\times(B\cap H)$ with this same $R$.
Thus intersection with $H$ is injective on those suffix flags and
preserves the selected nonzero coefficient. The diagram depicts a
flag in that support, not an arbitrary flag in the ambient poset.
When $c=0$, the initial flag is empty, $R=1$, and the lemma uses the
original cycle in $H$.}
\label{fig:fixed-complement}
\end{figure}
